\section{Orquestación multi-Agent: de la libertad estructural a la disciplina de ingeniería}

\subsection{Dimensiones de la orquestación}
El curso ofrece dimensiones de diseño claras para la orquestación multi-Agent: estático/dinámico, centralizado/descentralizado, contexto compartido/aislado, automatizado/con intervención humana. Estas dimensiones determinan el límite superior del desempeño del sistema y el límite inferior de su estabilidad.

\begin{figure}[H]
\centering
\includegraphics[width=0.93\textwidth]{frame-04-orchestration.jpg}
\caption{Dimensiones de diseño de Multi-Agent Orchestration}
\end{figure}
\noindent\footnotesize Evidencia visual: los subtítulos se ubican aproximadamente en 00:06:02--00:07:40, cuando el orador aborda la resolución de tareas complejas y la organización de trayectorias.\normalsize

\begin{importantbox}{Multi-Agent no es «múltiples ventanas en paralelo», sino «diseño de interfaces de tarea»}
Si no hay límites de rol, límites de contexto y límites de aceptación, aumentar el número de Agents por lo regular solo amplifica los conflictos y el trabajo duplicado.  
El paralelismo verdaderamente eficaz proviene de contratos claros de entrada y salida, así como de formatos claros para los productos intermedios.
\end{importantbox}

\subsection{Patrones de diseño: Conversation / Tool Use / Planning / Memory}
\begin{figure}[H]
\centering
\includegraphics[width=0.93\textwidth]{frame-05-patterns.jpg}
\caption{Agentic Design Patterns: colaboración entre conversación, herramientas, planificación y memoria}
\end{figure}
\noindent\footnotesize Evidencia visual: los subtítulos se ubican aproximadamente en 00:07:50--00:09:10, cuando el orador analiza sucesivamente patrones como reflection, plan y fact-check.\normalsize

\begin{knowledgebox}{Cómo los cuatro tipos de patrones forman un sistema mínimo viable}
\begin{itemize}
    \item Conversation: sostiene la colaboración entre roles y la transmisión de contexto;
    \item Tool Use: traduce las decisiones lingüísticas en acciones ejecutables;
    \item Planning: descompone tareas largas en subobjetivos verificables;
    \item Memory: conserva y comprime el estado clave a través de los turnos.
\end{itemize}
\end{knowledgebox}

\begin{warningbox}{Limitarse a «apilar patrones» hace que la complejidad del sistema se salga de control}
Los patrones no son mejores cuantos más se usen. Cada mecanismo que se agrega debe responder dos preguntas:  
¿qué tipo de tasa de fallos reduce? ¿qué tipo de nuevo modo de falla introduce?
\end{warningbox}

\subsection{Resumen de la sección}
El núcleo de la ingeniería multi-Agent no es la «orquestación creativa», sino el «diseño de restricciones». Solo cuando las interfaces son claras, el estado es verificable y la reversión es ejecutable, la orquestación compleja puede implementarse de manera estable.

